% Add in the preamble: \usepackage{graphicx}
% Russian captions require an existing Russian-language setup.
% Upload the seven PDFs to figures/remizov/ in Overleaf.
% Insert this fragment with \input{figures/remizov/figures.tex}.

\begin{figure}[htbp]
  \centering
  \includegraphics[width=0.92\linewidth]{figures/remizov/AMV1_F0000_S00_coefficients.pdf}
  \caption{Коэффициенты исходного уравнения AMV1\_F0000\_S00 на $[-1,1]$.}
  \label{fig:amv1_f0000_s00-coefficients}
\end{figure}

\begin{figure}[htbp]
  \centering
  \includegraphics[width=0.92\linewidth]{figures/remizov/AMV1_F0000_S00_solutions.pdf}
  \caption{Точное и сохранённое численное решения AMV1\_F0000\_S00. Маркеры обозначают 33 расчётных узла; штрихи соединяют соседние узлы.}
  \label{fig:amv1_f0000_s00-solutions}
\end{figure}

\begin{figure}[htbp]
  \centering
  \includegraphics[width=0.92\linewidth]{figures/remizov/AMV1_F0000_S00_absolute_error.pdf}
  \caption{Абсолютная ошибка AMV1\_F0000\_S00 в 33 сохранённых узлах. Пунктир обозначает порог $0.01$.}
  \label{fig:amv1_f0000_s00-absolute-error}
\end{figure}

\begin{figure}[htbp]
  \centering
  \includegraphics[width=0.92\linewidth]{figures/remizov/AMV1_F0141_S00_coefficients.pdf}
  \caption{Коэффициенты исходного уравнения AMV1\_F0141\_S00 на $[-1,1]$.}
  \label{fig:amv1_f0141_s00-coefficients}
\end{figure}

\begin{figure}[htbp]
  \centering
  \includegraphics[width=0.92\linewidth]{figures/remizov/AMV1_F0141_S00_solutions.pdf}
  \caption{Точное и сохранённое численное решения AMV1\_F0141\_S00. Маркеры обозначают 33 расчётных узла; штрихи соединяют соседние узлы. На вставке точное решение показано в собственном масштабе.}
  \label{fig:amv1_f0141_s00-solutions}
\end{figure}

\begin{figure}[htbp]
  \centering
  \includegraphics[width=0.92\linewidth]{figures/remizov/AMV1_F0141_S00_absolute_error.pdf}
  \caption{Абсолютная ошибка AMV1\_F0141\_S00 в 33 сохранённых узлах. Пунктир обозначает порог $0.01$.}
  \label{fig:amv1_f0141_s00-absolute-error}
\end{figure}

\begin{figure}[htbp]
  \centering
  \includegraphics[width=0.92\linewidth]{figures/remizov/AMV1_F0141_S00_absolute_error_log.pdf}
  \caption{Абсолютная ошибка AMV1\_F0141\_S00 в 33 сохранённых узлах; логарифмическая шкала по вертикали. Пунктир обозначает порог $0.01$.}
  \label{fig:amv1_f0141_s00-absolute-error-log}
\end{figure}
